%% ma: L12-B08-0120
%% lop: 12
%% chuong: 2
%% bai: 8
%% dang: 12.8.44
%% loai: TN4
%% muc: TH
%% dap_an: "B"
%% nguon: "CĐ2.3.pdf giáo viên gửi 27/10/2026 (M07, Lê Minh Khôi – Nguyễn Anh Khôi, Chương II Chủ đề 3 Biểu thức tọa độ của các phép toán vectơ), Câu 32"
%% kien_thuc: "Bài 8 (hai vectơ cùng phương, ba điểm thẳng hàng)"
%% trang_thai: cho-duyet
%% ly_do: "Dạng mới 12.8.44 đề xuất, chờ giáo viên duyệt"
\begin{ex}
Trong không gian $Oxyz$, cho ba điểm $A(2;-1;3)$, $B(-10;5;3)$ và $M(2m-1;2;n+2)$ (với $m$, $n$ là tham số). Để $A$, $B$, $M$ thẳng hàng thì giá trị của $m$, $n$ là
\choice{$m=1$, $n=\dfrac32$}{\True $m=-\dfrac32$, $n=1$}{$m=-1$, $n=-\dfrac32$}{$m=\dfrac23$, $n=\dfrac32$}
\loigiai{$\overrightarrow{AB}=(-12;6;0)$, $\overrightarrow{AM}=(2m-3;3;n-1)=t\overrightarrow{AB}\Rightarrow t=\dfrac12$.

$2m-3=-6\Rightarrow m=-\dfrac32$; $n-1=0\Rightarrow n=1$.}
\end{ex}
